\documentclass[a4paper,11pt]{article}
\hyphenation{acro-nym}
\usepackage{url}
\usepackage{xspace}
\usepackage[T1]{fontenc}
\usepackage{lmodern} % slightly better modern font
\usepackage{geometry}
\geometry{
    left=.5in,
    right=.5in,
    top=.5in,
    bottom=.5in
}

%\setlength{\parindent}{0pt}    % turn off indentation
\setlength{\parskip}{1em}   % how much to indent
\begin{document}
\title{Example 3}
\author{My name}
\date{January 5, 2011}
\maketitle
\section{What's this?}
This is our
second document.

It contains two paragraphs. The first line of a paragraph will be
indented, but not when it follows a heading.
% Here's a commen% Here's a comment.


50\% of \$100 makes \$50.
More special symbols are \&, \_, \{ and \}.





Text can be \emph{emphasized}.
Besides being \textit{italic} words could be \textbf{bold},
\textsl{slanted} or typeset in \textsc{Small Caps}.
Such commands can be \textit{\textbf{nested}}.
\emph{See how \emph{emphasizing} looks when nested.}



The best place for downloading LaTeX related software is CTAN.
Its address is \texttt{http://www.ctan.org}.





\section{\sffamily\LaTeX\ resources in the internet}
The best place for downloading LaTeX related software is CTAN.
Its address is \ttfamily http://www.ctan.org\rmfamily.


{\sffamily Text can be {\em emphasized}.  Besides being {\itshape italic} words could be {\bfseries bold}, {\slshape slanted} or typeset in {\scshape Small Caps}.  

Such commands can be {\itshape\bfseries nested}.}

{\em See ho {\em \bfseries emphasizing} looks when nested.}


%\noindent suppresses paragraph indentation
\noindent \tiny We \scriptsize start \footnotesize \small small \normalsize and \large big \Large Large  \LARGE and bigger \huge and \Huge gigantic!


\normalsize

% rare to declare fonts like this

\section{using environments}
\begin{huge}
\bfseries % this ends with the environment
A small example
%we ended paragraph before font size change

\end{huge}
\bigskip
This is just another illustrative example
\section{creating my own commands}

%\usepackage{xspace}
%\newcommand{\TUG}{TeX Users Group} %normally declare this in the preamble
%\newcommand{\TUG}{{\scshape \TeX\ Users Group}\xspace} %normally declare this in the preamble
\newcommand{\TUG}{\textsc{ \TeX\ Users Group}\xspace} %normally declare this in the preamble

\subsection{ The \TUG}
% we have to use the space
The \TUG is an organization for people who are interested in \TeX\ or \LaTeX
\subsection{commands with arguments}
%\newcommand{\keyword}[1]{\textbf{#1}}
\newcommand{\keyword}[2][\bfseries]{{#1#2}}

\keyword{Grouping} by curly braces limits the \keyword[\itshape]{scope} of \keyword{declarations}

\newcommand{\site}[2][https]{\textit{\texttt{#1://#2}}}

\site[ftp]{linux.org}


\url{www.nick.com}

%\parbox[alignment][height][inner alignment]{width}{text}
%inner aligntment is c,t,b,s center, top bottom stretch
% alignment is t or b

\subsection{parbox}

for small text boxes

\parbox{3cm}{TUG is an acronym.  It means \TeX\ Users group}


text line
\quad\parbox[b]{1.8cm}{this parbox is aligned at its bottom line}
\quad\parbox{1.5cm}{center-aligned parbox}
\quad\parbox[t]{2cm}{another parbox aligned at its top line}

\subsection{minipage}


\begin{minipage}{5cm}
  TUG is an acronym. It means \TeX\ Users Group. \footnote{my first footnote in a while}
\end{minipage}

That was an environment, not a command. But they take optional arguments and arguments like a command.

\subsection{manual hyphenation}
% put tht in preamble
%\hyphenation{acro-nym}
\parbox{3cm}{TUG is an acronym.  It means \TeX\ Users group}
\quad


\newcommand{\packex}[1]{\texttt{\textbackslash#1}}

can go \texttt{\textbackslash\{indivisible\}} that has no break points so it will not hyphenate

or \texttt{\textbackslash\{mbox\}} which disables all breaks

there is also the \texttt{hyphenat} package. 

\packex{ladeeda}








\end{document}

